\section{System Initialization and Boot Flow}
\label{sec:vayu_boot}

The system initialization and boot flow in Vayu defines the sequence of operations that transition the system from power-on to a fully operational state. This process ensures that hardware, system services, and application tasks are correctly initialized before entering normal operation.

\subsection{Initialization Sequence}

System startup begins with low-level hardware configuration, followed by initialization of system services and application components. The sequence is structured to ensure that all dependencies are satisfied before tasks begin execution.

The initialization flow proceeds as follows:

\begin{enumerate}
    \item \textbf{Clock Configuration:}  
    The system clock is configured using the phase-locked loop (PLL) to achieve the required operating frequency. This ensures consistent timing for all subsystems.

    \item \textbf{Core System Initialization:}  
    VAIOS is initialized using a configuration structure that defines system behavior, including clock handling and peripheral setup.

    \item \textbf{Peripheral Initialization:}  
    Hardware interfaces such as I2C and UART are initialized through NavHAL. This includes setting up DMA for sensor acquisition and communication channels.

    \item \textbf{Subsystem Initialization:}  
    Core subsystems including logging, system state management, sensor interfaces, and communication buffers are initialized.

    \item \textbf{Task Creation:}  
    System and application tasks are created and registered with the scheduler. These include control tasks, sensor acquisition tasks, communication handlers, and auxiliary tasks.

    \item \textbf{Timer Setup:}  
    Periodic timer callbacks are configured to support time-based operations and high-frequency timing requirements.

    \item \textbf{Scheduler Start:}  
    The VAIOS scheduler is started, after which all task execution is managed by the RTOS.
\end{enumerate}

\subsection{Boot Task and System Checks}

A dedicated boot task performs system validation checks during startup. This task is responsible for verifying that critical components are functioning correctly before enabling normal operation.

The boot checks include:
\begin{itemize}
    \item System clock validation
    \item Storage subsystem availability
    \item Basic system readiness verification
\end{itemize}

The results of these checks are accumulated into a boot status register, which is used to determine the final system state.

If all required checks pass, the system transitions to a standby state, indicating readiness for operation. If any check fails, the system enters a failsafe state to prevent unsafe behavior.

\subsection{State Transition}

The system state progresses through a well-defined sequence:

\[
\text{INIT} \rightarrow \text{STANDBY} \rightarrow \text{ARMED}
\]

During initialization, the system remains in the \textit{INIT} state until all boot checks are completed. Once validated, it transitions to \textit{STANDBY}, where it awaits user input. The transition to the \textit{ARMED} state is controlled by external input (e.g., RC commands).

\subsection{Task Activation}

After the scheduler starts, all created tasks begin execution based on their assigned priorities. High-priority tasks such as control loops and sensor acquisition become active immediately, while lower-priority tasks handle communication, logging, and background operations.

The boot task terminates itself after completing initialization, ensuring that it does not consume system resources during normal operation.

\subsection{System Behavior}

The structured initialization flow ensures that all components are properly configured before entering real-time execution. By separating initialization, validation, and runtime operation, the system maintains robustness and avoids undefined behavior during startup.
